\documentclass[../main.tex]{subfiles}

\begin{document}

\ifSubfilesClassLoaded{

    \tableofcontents
    
}{}

\chapter{ Law of Large Number }
\section{Independence}
\begin{definition}{Independence}{}
    \begin{itemize}
        \item  Two events $A$ and $B$ are \dfntxt{independent} if $P(A \cap B) = P(A)P(B)$.
        \item Two random variables $X$ and $Y$ are \dfntxt{independent} if for all $C, D \in \mathcal{R}$,
\[
\begin{aligned}
P(X \in C, Y \in D) = P(X \in C)P(Y \in D)
\end{aligned}
\]that is, the events $A = \{X \in C\}$ and $B = \{Y \in D\}$ are independent.
\item Two $\sigma$-fields $\mathcal{F}$ and $\mathcal{G}$ are \dfntxt{independent} if for all $A \in \mathcal{F}$ and $B \in \mathcal{G}$ the events $A$ and $B$ are independent.
    \end{itemize}
\end{definition}


\begin{exercise}{}{}
\begin{enumerate}
    \item Show that if $A$ and $B$ are independent, then so are $A^c$ and $B$, $A$ and $B^c$, and $A^c$ and $B^c$.
    \item Conclude that events $A$ and $B$ are independent if and only if their indicator random variables $1_A$ and $1_B$ are independent.
\end{enumerate}
\end{exercise}
\begin{proof}
    \begin{enumerate}
        \item If \(  A  \) and \(  B  \) are independent, then \[
        P\left( A\cap B \right)= P\left( A \right)P\left( B \right)   
        \]   Note that \[
        P\left( A\cap B \right)+ P\left( A^{c}\cap B \right)= P\left( B \right)   
        \] We have \[
        P\left( A^{c}\cap B \right)= P\left( B \right)\left( 1-P\left( A \right)  \right)= P\left( B \right)P\left( A^{c} \right)     
        \]  Similarly, \[
        P\left( A\cap B^{c} \right)= P\left( A \right)P\left( B^{c} \right)   
        \] \[
        P\left( A^{c}\cap B^{c} \right)= P\left( B^{c} \right)-P\left( A\cap B^{c} \right)= P\left( B^{c} \right)- P\left( A \right)P\left( B^{c} \right)= P\left( A^{c} \right)P\left( B^{c} \right)        
        \]
        \item  \[
         \sigma \left( 1_{A} \right)= \left\{ \varnothing,A, A^{c}, \Omega  \right\} ,\quad  \sigma \left( 1_{B} \right)= \left\{ \varnothing,B,B^{c}, \Omega  \right\} 
        \] If \(   \sigma \left( 1_{A} \right)   \) and \(   \sigma \left( 1_{B} \right)   \) are independent, it is clear that \(  A  \) and \(  B  \) are as well.  Conversly, only need to show that  \(  \varnothing  \) and \(  A  \) are independent, the remaining follows from 1. and symmetry.        We have \[
        P\left( \varnothing \cap A \right) = 0=  P\left( \varnothing \right)P\left( A \right)   
        \]
    \end{enumerate}
    
\end{proof}

\begin{definition}{}{}
    \begin{itemize}
        \item $\sigma$-fields $\mathcal{F}_1, \mathcal{F}_2, \dots, \mathcal{F}_n$ are \dfntxt{independent} if whenever $A_i \in \mathcal{F}_i$ for $i = 1, \dots, n$, we have
\[
P \left(\bigcap_{i=1}^n A_i\right) = \prod_{i=1}^n P(A_i)
\]
\item 
Random variables $X_1, \dots, X_n$ are \dfntxt{independent} if whenever $B_i \in \mathcal{R}$ for $i = 1, \dots, n$ we have
\[
P \left(\bigcap_{i=1}^n \{X_i \in B_i\}\right) = \prod_{i=1}^n P(X_i \in B_i)
\]
\item Sets $A_1, \dots, A_n$ are \dfntxt{independent} if whenever $I \subset \{1, \dots, n\}$ we have
\[
P \left(\bigcap_{i \in I} A_i\right) = \prod_{i \in I} P(A_i)
\]
    \end{itemize}
\end{definition}

\begin{remark}
    The last definition looks slightly different from the former two, but it is consistent with the independence of the indicator variables \(  1_{A_{i}},1\le i\le n  \).  
\end{remark}
\begin{exercise}{}{}
Let $A_1, A_2, \dots, A_n$ be independent. Show
\begin{enumerate}
    \item $A_1^c, A_2, \dots, A_n$ are independent;
    \item $1_{A_1}, \dots, 1_{A_n}$ are independent.
\end{enumerate}
\end{exercise}
\begin{proof}
  \begin{enumerate}
    \item   If \(  1\not \in I  \), then it is obvious that  \[
    P\left( \cap _{i\in I}A_{i} \right)= \prod _{i\in I}P\left( A_{i} \right).   
    \]  If \(  1 \in I  \), note that    \[
   \begin{aligned}
    P\left( \cap _{i \in I\setminus \left\{ 1 \right\}}A_{i}\cap  A_{1} ^{c}\right)&=  P\left( \cap _{i \in I\setminus \left\{ 1 \right\}} A_{i}\right)-  P\left( \cap _{i \in I\setminus \left\{ 1 \right\}}A_{i}\cap A_{1} \right)  \\ 
     &=  \prod _{i\in I\setminus \left\{ 1 \right\}}P\left( A_{i} \right)-\prod _{ i\in I}P\left( A_{i} \right)   \\ 
      &=  \left( 1-P\left( A_{1} \right)  \right)\prod _{ i \in I\setminus \left\{ 1 \right\}} P\left( A_{i} \right)\\ 
       &=  P\left( A_{1}^{c} \right)\prod _{i \in I\setminus \left\{ 1 \right\}}P\left( A_{i} \right)   
   \end{aligned}
    \]
    \item  \[
     \sigma \left( 1_{A_{i}} \right) = \left\{ \varnothing,A_{i}^{c},A_{i}, \Omega  \right\}
    \]  Take \(  B_{i}\in  \sigma \left( 1_{A_{i}} \right), i=  1,\cdots,n    \).  \(  B_{i}= \varnothing  \) for some \(  i  \)   , then both \(  P\left( \bigcap_{i= 1}^{n}B_{i}  \right)   \) and \(  \prod _{i= 1}^{n}P\left( B_{i} \right)   \) are zero.   The ramainings follows from 1. 
  \end{enumerate}
  
\end{proof}

\begin{example}{}{}
Let $X_1, X_2, X_3$ be independent random variables with
\[
P(X_i = 0) = P(X_i = 1) = 1/2
\]
Let $A_1 = \{X_2 = X_3\}$, $A_2 = \{X_3 = X_1\}$ and $A_3 = \{X_1 = X_2\}$. These events are
pairwise independent since if $i \neq j$, then
\[
P(A_i \cap A_j) = P(X_1 = X_2 = X_3) = 1/4 = P(A_i)P(A_j)
\]
but they are not independent since
\[
P(A_1 \cap A_2 \cap A_3) = 1/4 \neq 1/8 = P(A_1)P(A_2)P(A_3)
\]
\end{example}
\subsection{Sufficient Conditions}

\begin{definition}{Independence between conllections}{}
    Collections of sets $\mathcal{A}_1, \mathcal{A}_2, \ldots, \mathcal{A}_n \subset \mathcal{F}$ are said to be \dfntxt{independent} if whenever $A_i \in \mathcal{A}_i$ and $I \subset \{1, \ldots, n\}$ we have $P(\bigcap_{i \in I} A_i) = \prod_{i \in I} P(A_i)$.
\end{definition}

\begin{lemma}{}{}
Without loss of generality we can suppose each $\mathcal{A}_{i}$ contains $\Omega$. In this case the condition is equivalent to
\[
P \left(\bigcap_{i=1}^n A_i\right) = \prod_{i=1}^n P(A_i) \text{ whenever } A_i \in \mathcal{A}_i
\]
since we can set $A_i = \Omega$ for $i \notin I$.
\end{lemma}

\begin{definition}{\(  \pi  \),\(   \lambda   \)- systems  }{}
    \begin{itemize}
        \item  We say that $\mathcal{A}$ is a \dfntxt{$\pi$-system} if it is closed under intersection, that is, if $A, B \in \mathcal{A}$ then $A \cap B \in \mathcal{A}$. 
        \item     
 We say that $\mathcal{L}$ is a \dfntxt{$\lambda$-system} if
\begin{enumerate}
    \item[(i)] $\Omega \in \mathcal{L}$.
    \item[(ii)] If $A, B \in \mathcal{L}$ and $A \subset B$, then $B - A \in \mathcal{L}$.
    \item[(iii)] If $A_n \in \mathcal{L}$ and $A_n \uparrow A$, then $A \in \mathcal{L}$.
\end{enumerate}
    \end{itemize}

\end{definition}
\begin{theorem}{$\pi - \lambda$ \dfntxt{Theorem}}{9-3-3}
If $\mathcal{P}$ is a $\pi$-system and $\mathcal{L}$ is a $\lambda$-system that contains $\mathcal{P}$, then $\sigma(\mathcal{P}) \subset \mathcal{L}$.
\end{theorem}
\begin{theorem}{}{8-3-1}
Suppose $\mathcal{A}_1, \mathcal{A}_2, \dots, \mathcal{A}_n$ are independent and each $\mathcal{A}_i$ is a $\pi$-system. Then $\sigma(\mathcal{A}_1), \sigma(\mathcal{A}_2), \dots, \sigma(\mathcal{A}_n)$ are independent.
\end{theorem}
\begin{proofsketch}
     对于任一取定的事件 \(  F \),  ``(取定任一集族中)与 \(  F  \)独立 ''的性质集是一个 \(   \lambda   \)-system. 由于 \(  \mathcal{A}_{1}  \)是 \(   \pi  \)-system,  于是 \(  \mathcal{A}_{1}  \)中与 \(  F  \)独立 的集合包含了 \(   \sigma \left( \mathcal{A}_{1} \right)   \), 通过任意的 \(  F  \), 传递归纳 集族的独立性一步步``升级''到 它们的生成 \(   \sigma   \)-代数上.   
\end{proofsketch}
\begin{remark}
    It is not easy to show that if \(  A,B \in \mathcal{L}  \) then \(  A\cap B\in \mathcal{L}  \) or \(  A\cup B\in \mathcal{L}  \), but it is easy to check that if \(  A,B\in \mathcal{L}  \) with \(  A\subseteq B  \), then \(  B-A\in \mathcal{L}  \).      
\end{remark}
\begin{proof}
    Let \( F\in \mathcal{P}\) be any event. We define \[
    \mathcal{L}: = \left\{ A\in \mathcal{A}_{1}: P\left( A\cap F \right)= P\left( A \right)P\left( F \right)    \right\}
    \]We check that \(  \mathcal{L}  \)   is a \(   \lambda   \)-system. (i) follows from  \(  P\left( A\cap  \Omega  \right)= P\left( A \right)\cdot 1= P\left( A \right)P\left(  \Omega  \right)      \). To show (ii), we take  \( A, B \in \mathcal{L}, B\supseteq A  \). Note that \[
    \left( B-A \right)\cap F=  \left( B\cap F \right)-\left( A\cap F \right),  
    \] we have \[
    \begin{aligned}
    P\left( \left( B\setminus A \right)\cap F  \right)&= P\left( \left( B\cap F \right)-\left( A\cap F \right)   \right)\\ 
     &=    P\left( B\cap F \right)-P\left( A\cap F \right)\\ 
      &=   P\left( B \right)P\left( F \right)-P\left( A \right)P\left( F \right)\\ 
       &= \left( P\left( B \right)-P\left( A \right)   \right)P\left( F \right)=  P\left( \left( B-A \right)  \right)P\left( F \right)       
    \end{aligned}
    \]   which shows that \(  \left( B-A \right)\in \mathcal{L}   \) 

    For (iii), take \(  B_{n}\in \mathcal{L}, B_{n}\uparrow B \). Then since \(  \left( B_{n}\cap F \right)\uparrow B\cap F   \), we have \[
    \begin{aligned}
    P\left( B\cap F \right)&= \lim_{n\to \infty} P\left( B_{n}\cap F \right)\\ 
     &= \lim_{n\to \infty}P\left( B_{n} \right)P\left( F \right)\\ 
      &=  P\left( B \right)P\left( F \right)      
    \end{aligned},
    \]which implies that \(  B \in \mathcal{L}  \). 
    Thus \(  \mathcal{L}  \) is a \(   \lambda   \)-system. It follows from \(   \pi- \lambda   \) Theorem that \(  \mathcal{L}\supseteq  \sigma \left( \mathcal{A}_{1} \right)   \)  For each \(  A_2,\cdots ,A_{n}  \), \(  A_{i}\in \mathcal{A}_{i},i= 2,3,\cdots ,n  \), by taking \(  F= \bigcap_{i= 2}^{n}A_{i}   \), we have \[
    P\left( A \cap A_2\cap \cdots \cap A_{n}\right)= P\left( A \right)P\left( \bigcap_{i= 2}^{n}A_{i}  \right)= P\left( A \right)P\left( A_2 \right)\cdots P\left( A_{n} \right),\quad \forall  A\in  \sigma \left( \mathcal{A}_{1} \right)       
    \]      We get that \[
\tag{*} \text{if }   \mathcal{A}_{1},\cdots ,\mathcal{A}_{n}  \text{ are independent }  \pi-\text{systems, then }  \sigma \left( \mathcal{A}_{1} \right),\mathcal{A}_{2},\cdots ,\mathcal{A}_{n}\text{ are as well}  
    \]  By applying \(  (*)  \) into \(  \mathcal{A}_{2},\cdots ,\mathcal{A}_{n},  \sigma \left( \mathcal{A}_{1} \right)   \), we get \(   \sigma \left( \mathcal{A}_{2} \right),\mathcal{A}_{3},\cdots ,\mathcal{A}_{n}, \sigma \left( \mathcal{A}_{1} \right)    \)    are independent. The proof completes after \(  n  \) iterations. 
\end{proof}
 
\begin{theorem}{}{9-3-2}
In order for $X_1, \dots, X_n$ to be independent, it is sufficient that for all $x_1, \dots, x_n \in (-\infty, \infty]$
\[ P(X_1 \le x_1, \dots, X_n \le x_n) = \prod_{i=1}^n P(X_i \le x_i) \]
\end{theorem}
\begin{proof}
    Let \(  \mathcal{A}_{i}  \) be the colletcion of the sets of the form \(\left\{ X_{i}\le x_{i} \right\}  \). Since \[
    \left\{ X_{i}\le x\right\}\cap \left\{ X_{i}\le y \right\}= \left\{ X_{i}\le x\wedge y \right\}
    \]\(  \mathcal{A}_{i}  \) is a \(   \pi  \)-system. Since we allowed \(  x_{i}= \infty  \), \(   \Omega \in \mathcal{A}_{i}  \). Exercise \ref{ex:9-3-1} implies that \(   \sigma \left( \mathcal{A}_{i} \right)=  \sigma \left( X_{i} \right)    \), so the result follows from   Theorem \ref{thm:9-3-2}.  
\end{proof}
\begin{theorem}{}{}
Suppose $\mathcal{F}_{i,j}$, $1 \leq i \leq n$, $1 \leq j \leq m(i)$ are independent \(   \pi  \)-systems   and let $\mathcal{G}_i = \sigma(\bigcup_j \mathcal{F}_{i,j})$. Then $\mathcal{G}_1, \dots, \mathcal{G}_n$ are independent.
\end{theorem}
\begin{note}
    对于独立的 \(   \pi  \)-systems, 可以把它们按任意分组方式, ``攒起来''做 \(   \sigma   \)-代数扩张, 得到新的独立系统.  
\end{note}
\begin{proof}
Let $\mathcal{A}_i$ be the collection of sets of the form $\bigcap_j A_{i,j}$ where $A_{i,j} \in \mathcal{F}_{i,j}$. $\mathcal{A}_i$ is a $\pi$-system that contains $\Omega$ and contains $\bigcup_j \mathcal{F}_{i,j}$, so Theorem \ref{thm:8-3-1} implies $\sigma(\mathcal{A}_i) = \mathcal{G}_i$ are independent.
\end{proof}
\begin{theorem}{}{}
If for $1 \le i \le n$, $1 \le j \le m(i)$, $X_{i,j}$ are independent and $f_i : \mathbf{R}^{m(i)} \to \mathbf{R}$ are measurable, then $f_i(X_{i,1}, \dots, X_{i,m(i)})$ are independent.
\end{theorem}
\begin{proof}
     \(   \sigma \left(X_{i,j} \right), 1\le i\le n,1\le j\le m\left( i \right)    \) are independent \(   \pi  \)- systems. Then \(  \mathcal{G}_{i} =  \sigma \left( \bigcup _{j}\left(  \sigma \left( X_{i,j} \right)  \right)  \right)   \) are independent. Take any Borel set \(  B \subseteq \mathbb{R}   \), then \[
  \begin{aligned}
     \left( f_{i}\left( X_{i,1},\cdots ,X_{i,m\left( i \right) } \right)  \right)^{-1} \left( B \right)&=  \left( X_{i,1},\cdots ,X_{i,m\left( i \right) } \right)^{-1} \left( f_{i}^{-1} \left( B \right)  \right)\\ 
      &\subseteq   \sigma \left( \bigcup _{j}\left(  \sigma \left( X_{i,j} \right)  \right)  \right)   
  \end{aligned}  
     \] Thus \(  f_{i}\left( X_{i,1},\cdots ,X_{i,m\left( i \right) } \right)   \) are independent.  
\end{proof}

\subsection{Independence, Distribution, and Expectation}

\begin{theorem}{}{}
Suppose $X_1, \dots, X_n$ are independent random variables and $X_i$ has distribution $\mu_i$. Then $(X_1, \dots, X_n)$ has distribution $\mu_1 \times \dots \times \mu_n$.
\end{theorem}
\begin{proof}
     \[
     \begin{aligned}
     P\left( \left(  X_1,\cdots,X_n  \right)\in A_1\times \cdots \times A_{n}  \right)&= P\left( X_1\in A_1,\cdots ,X_{n}\in A_{n} \right)\\ 
      &=    \prod _{i= 1}^{n}P\left( X_{i}\in A_{i} \right)\\ 
       &=  \prod _{i= 1}^{n}\mu _{i}\left( A_{i} \right)\\ 
        &= \left( \mu _1 \times \cdots \times \mu _{n} \right)\left( A_1\times \cdots \times A_{n} \right)    
     \end{aligned}
     \]The above shows that the distribution of \(  (X_1,\cdots ,X_{n})  \) and the measure \(  \mu _1 \times \cdots \times \mu _{n}  \) agree on the sets of the form \(  A_1\times \cdots \times A_{n}  \), which forms a \(   \pi  \)-system  generates \(  \mathcal{R}^{n}  \).   Since the collection of the sets on which two measures here agree forms a \(  \mathcal{L}  \)-system. Then it follows from Theorem \ref{thm:9-3-3} that \(  \mu _1 \times \cdots \times \mu _{n}  \) and \(  \left( X_1,\cdots ,X_{n} \right)   \) agree on \(  \mathcal{R}^{n}  \).      
\end{proof}
\begin{theorem}{}{}
Suppose $X$ and $Y$ are independent and have distributions $\mu$ and $\nu$. If $h : \mathbf{R}^2 \to \mathbf{R}$ is a measurable function with $h \ge 0$ or $E|h(X, Y)| < \infty$, then
\[
\begin{aligned}
Eh(X, Y) = \iint h(x, y) \mu(dx) \nu(dy)
\end{aligned}
\]
In particular, if $h(x, y) = f(x)g(y)$ where $f, g : \mathbf{R} \to \mathbf{R}$ are measurable functions with $f, g \ge 0$ or $E|f(X)|$ and $E|g(Y)| < \infty$, then
\[
\begin{aligned}
Ef(X)g(Y) = Ef(X) \cdot Eg(Y)
\end{aligned}
\]
\end{theorem}

\begin{proof}
    From the Changing of Variables and the Fubini's Theorem, we have \[
    \begin{aligned}
    Eh\left( X,Y \right)&= \int h \,\mathrm{d} \left( \mu \times \nu  \right)= \iint h\left( x,y \right)\mu \left( \,\mathrm{d} x \right)\nu \left( \,\mathrm{d} y \right)      
    \end{aligned}
    \]

    To show the second part, we first assmu that \(  f,g \ge 0  \), then \[
    \begin{aligned}
     Ef\left( X \right)g\left( Y \right)&= \iint f\left( x \right)g\left( y \right)\mu \left( \,\mathrm{d} x \right)\nu \left( \,\mathrm{d} y \right)\\ 
      &= \int g\left( y \right) \int f\left( x \right)\mu \left( \,\mathrm{d} x \right)\nu \left( \,\mathrm{d} y \right)\\ 
       &=  \int g\left( y \right)Ef\left( X \right)          \nu \left( \,\mathrm{d} y \right)\\ 
        &=  Ef\left( X \right)Eg\left( Y \right)   
    \end{aligned}
    \] 
    For another assumption, we apply the above result to the nonnegative function \(  \left| f \right|, \left| g \right|    \), to get \(  E\left| h\left( X,Y \right)  \right|< \infty   \)  , and we repeat the last argument to prove the desired result.
\end{proof}
\begin{theorem}{}{}
If $X_1, \dots, X_n$ are independent and have
\begin{enumerate}
    \item[(a)] $X_i \ge 0$ for all $i$, or
    \item[(b)] $E|X_i| < \infty$ for all $i$, then
\end{enumerate}
\[
E\left(\prod_{i=1}^n X_i\right) = \prod_{i=1}^n EX_i
\]
that is, the expectation on the left exists and has the value given on the right.
\end{theorem}


\subsection{Sums of Independent Random Variables}

\begin{theorem}{}{}
If $X$ and $Y$ are independent, $F(x) = P(X \leq x)$, and $G(y) = P(Y \leq y)$, then
\[
P(X + Y \leq z) = \int F(z - y) dG(y)
\]
The integral on the right-hand side is called the \dfntxt{convolution} of $F$ and $G$ and is denoted $F * G(z)$. The meaning of $dG(y)$ will be explained in the proof.
\end{theorem}

\begin{proof}
    \[
   \begin{aligned}
    P\left( X+ Y\le z \right)&=  \int 1_{\left( x+ y \le z\right) }\,\mathrm{d} \left( \mu \times \nu  \right)\\ 
     &=   \int \int 1_{\left( x+ y\le z \right) } \,\mathrm{d} \mu  \,\mathrm{d} \nu \\ 
      &=  \int F\left( z-y \right)\,\mathrm{d} \nu \\ 
       &=  \int F\left( z-y \right) \,\mathrm{d} G\left( y \right)   
   \end{aligned} 
    \]where \(  \,\mathrm{d} G\left( y \right)   \) is just a notation for \(  \,\mathrm{d} \nu   \).  We apply this notation motivated from the Lebesgue-Stieljes integration.  
\end{proof}

\begin{theorem}{}{}
Suppose that $X$ with density $f$ and $Y$ with distribution function $G$ are independent. Then $X + Y$ has density
\[
h(x) = \int f(x - y)dG(y)
\]
When $Y$ has density $g$, the last formula can be written as
\[
h(x) = \int f(x - y) g(y)dy
\]
\end{theorem}

\begin{proof}
     \[
     F\left( x \right)=  \int_{-\infty}^{x}f\left( x \right)\,\mathrm{d} x  
     \] Since everything is nonnegative, from Fubini's theorem, we get\[
  \begin{aligned}
     P\left( X+ Y\le z \right)&=  \int \int_{-\infty}^{z-y}f\left( x \right)\,\mathrm{d} x \,\mathrm{d} G\left( y \right)    \\ 
      &=  \int \int_{-\infty}^{z}f\left( x- y \right)\,\mathrm{d} x\,\mathrm{d} G\left( y \right)\\ 
       &= \int_{-\infty}^{z}\int f\left( x-y \right)\,\mathrm{d} G\left( y \right)\,\mathrm{d} x  \\ 
        &=  \int_{-\infty}^{z} h\left( x \right)\,\mathrm{d} x   
  \end{aligned}
     \] which shows that \(  X+ Y   \) has density \(  h\left( x \right)   \). 
     
     When \(  Y  \) has density \(  g  \), then \[
    \int f\left( x-y \right)\,\mathrm{d}  G\left( y \right)=   \int f\left( x-y \right) g\left( y \right)\,\mathrm{d} y 
     \] 
\end{proof}
\begin{example}{The gamma density}{}
The \dfntxt{gamma density} with parameters $\alpha$ and $\lambda$ is given by
\[
f(x) = \begin{cases}
\lambda^\alpha x^{\alpha-1}e^{-\lambda x} / \Gamma(\alpha) & \text{for } x \geq 0 \\
0 & \text{for } x < 0
\end{cases}
\]
where $\Gamma(\alpha) = \int_0^\infty x^{\alpha-1}e^{-x}dx$.
\end{example}

\begin{theorem}{}{}
    If \(  X = \mathrm{gamma}\left( \alpha , \lambda  \right)   \) and \(  Y= \operatorname{gamma}\left( \beta , \lambda  \right)   \) are independent, then \(  X+ Y  \)    is \(  \operatorname{gamma}\left( \alpha + \beta , \lambda  \right)   \). Consequently if \(   X_1,\cdots,X_n   \) are independent \(  \operatorname{exponential}\left(  \lambda  \right)    \) r.v's, then \(  X_1+ \cdots + X_{n}  \) has a \(  \operatorname{gamma}\left( n, \lambda  \right)   \) distribution.      
\end{theorem}

\begin{example}{Normal distribution}{}
In Example 1.6.2, we introduced the normal
density with mean $\mu$ and variance $a$,
\[
(2\pi a)^{-1/2} \exp(-(x - \mu)^2/2a).
\]
\end{example}

\begin{theorem}{}{}
If $X = \text{normal}(\mu, a)$ and $Y = \text{normal}(\nu, b)$ are independent, then $X + Y = \text{normal}(\mu + \nu, a + b)$.
\end{theorem}

\section{Weak Laws of Large Numbers}
\end{document}